\section{The four-copy transport operator}\label{sec:transport}

We now explain the rank jump behind \((6.4)\).  Fix the six copy
permutations
\[
\boldsymbol\sigma=(1,(03)(12),(23),(02)(13),(01)(23),(021)).
\]
Let \(V\) be the three-dimensional message space and let
\(g_i:V\to k\) be the six coordinate covectors.  Form a bipartite
four-sheeted cover with vertices \(L_a,R_b\), \(0\leq a,b<4\), and a
color-\(i\) edge from \(L_a\) to \(R_{\sigma_i(a)}\).  Put \(V\) at each
vertex and \(k\) at each edge, using \(g_{p(i)}\) as both restriction
maps for a party assignment \(p\in S_6\).  A global section is a tuple
\((x_a,y_b)\in V^8\) satisfying
\[
 g_{p(i)}(x_a-y_{\sigma_i(a)})=0                       \tag{7.1}
\]
on all \(24\) edges.  The three-dimensional constant-section space is
always present.  Gauging it away leaves the \(24\)-by-\(21\) matching
matrix in \((6.3)\).

Let \(U=k^4/k(1,1,1,1)\), and let \(\rho_i\) be the action of
\(\sigma_i\) on \(U\).  For each \(p\), the MDS property makes
\(g_{p(3)},g_{p(4)},g_{p(5)}\) a basis of \(V^*\).  Define \(Q_p\) by
\[
 g_{p(i)}=\sum_{k=0}^2(Q_p)_{ik}g_{p(3+k)}
 \qquad(0\leq i<3).                                   \tag{7.2}
\]
Thus \(Q_p\) is the left block obtained by systematizing the
party-reordered generator matrix on its last three columns; no pivot can
vanish because every three coordinate covectors are independent.  For
the identity assignment,
\[
Q_1=\begin{pmatrix}a&b&c\\a&c&b\\d&d&d\end{pmatrix},
\quad
\begin{array}{ll}
a=-2(t-1)^2,&b=-(t^2-t+1),\\
c=-(t^2-3t+1),&d=-2(t-1).
\end{array}
\]
If \(Y_k\in U\) are the three free copy vectors, then \((7.1)\) reduces
to the \(9\)-by-\(9\) block transport operator
\[
 T_p=\bigl((Q_p)_{ik}(\rho_i-\rho_{3+k})\bigr)_{i,k=0}^2.
                                                               \tag{7.3}
\]
The dimension identity
\[
 21-\rank M_{\pi\boldsymbol\sigma}(G(t))=9-\rank T_p    \tag{7.4}
\]
holds over the field \(k\), where \(p=\pi^{-1}\) in the convention of
Section~\ref{sec:invariants}.  It follows by quotienting the
three-dimensional constant-section kernel and then eliminating the three
free systematic coordinates.  Both steps are invertible on the six-arc
locus.

\begin{theorem}[Transport divisor]\label{thm:transport-divisor}
For a parameter on the admitted odd non-GRS pencil, there exists a party
assignment \(p\in S_6\) for which the sheaf \((7.1)\) has a nonconstant
section if and only if
\[
 (B^2-2A^2)(9B^2-4A^2)=0,
\]
or equivalently
\[
 (z-2)(9z-4)=0.                                       \tag{7.5}
\]
At a generic point above \(z=2\), exactly \(96\) party assignments have
rank \(20\); at a generic point above \(z=4/9\), exactly \(192\) do.
Every other assignment has rank \(21\).
\end{theorem}

\begin{proof}
For the identity systematic chart, the coefficients satisfy
\[
 a=b+c,\qquad d^2=-2a,\qquad A=a(c-b),\qquad B=bc.
\]
Common copy relabelling, a change of basis in \(V\), and rescaling a
coordinate covector act on \(T_p\) by invertible row and column
operations.  Modulo these operations, the \(720\) party assignments
split into boundary types and three nonboundary relative-cover types.
Here the type records the six relative copy permutations
\(\rho_j^{-1}\rho_i\), equivalently the even-cycle pattern in each pair
of colored matchings.  Expanding a representative determinant by signed
cycle covers gives
\[
\begin{split}
P_-&=AB(3B-2A),\\
P_+&=AB(3B+2A),\\
P_2&=A(B^2-2A^2).
\end{split}
\]
With integral bases the determinants are respectively
\[
64d^3AB(3B-2A),\quad
64d^3AB(3B+2A),\quad
64d^3A(B^2-2A^2).
\]
All omitted factors are units on the admitted odd pencil.  Every
assignment outside these three types has a determinant that is a unit
there.  The signed cycle-cover ledger, the exhaustion of the \(720\)
assignments, and a fraction-free determinant path are recorded in the
paper-local certificate.  The two signed types lie above
\(w=B/A=\pm2/3\), hence \(z=4/9\); the axial type lies above \(z=2\).
Fraction-free elimination also gives a nonzero \(8\)-minor on each
component, so the generic kernel is exactly one-dimensional.

The row/column equivalences act on assignments from the left by copy
symmetries and from the right by color stabilizers.  The axial support is
one double coset with groups \(S_4\times C_2\) and \(D_8\times C_2\)
and intersection \(C_2^3\).  Its size is
\[
\frac{|S_4\times C_2|\,|D_8\times C_2|}{|C_2^3|}
 =\frac{48\cdot16}{8}=96.
\]
Each signed support is one double coset with groups \(S_4\) and
\(S_4\times C_2\) and intersection \(S_3\), hence has size
\[
\frac{|S_4|\,|S_4\times C_2|}{|S_3|}
 =\frac{24\cdot48}{6}=192.
\]
The signed types are distinguished by the sign of \(w\), and the axial
cycle pattern differs from both, so the three supports are disjoint.
The orbit certificate supplies explicit representatives, stabilizers,
intersections, and the complete support sets.  Equation \((7.4)\) turns
the one-dimensional transport kernels into rank \(20\) of the matching
matrix and completes the proof.
\end{proof}

The exceptional arithmetic belongs to this detector, not to
Theorem~\ref{thm:lu-lc-rigidity}.  The exact boundary is:
\[
\begin{array}{c|l}
\operatorname{char}k&\text{effect on the transport divisor}\\ \hline
2&\text{outside the odd-field pencil}\\
3&\text{both admitted detector components disappear}\\
5&z=4/9\text{ lies on the excluded GRS boundary; }z=2\text{ remains}\\
7&\text{the two components merge; the rank histogram is }20^{288}21^{432}\\
11,13,41&\text{pullback ramification only; no further rank drop.}
\end{array}                                             \tag{7.6}
\]
At \(t=-1\), \(A=16\), \(B=15\), and the nonboundary exceptional primes
are already visible in
\[
B^2-2A^2=-7\cdot41,\qquad
9B^2-4A^2=7\cdot11\cdot13.
\]
In characteristic \(3\), the axial factor becomes \(G^2\) and the signed
factors lie on \(t(t-1)=0\), so no admitted component remains.  In
characteristic \(5\), \(9B^2-4A^2\) is a unit times \(G^2\), removing
only the signed component.  In characteristic \(7\),
\(9B^2-4A^2=2(B^2-2A^2)\); the disjoint \(96\)- and \(192\)-element
supports merge, and exact \(8\)-minors exclude a larger kernel.
At \(11,13,41\) the indicated pullback component has a double root at
\(t=-1\), but the same minors keep the rank drop equal to one.

The divisor records where one scalar contraction changes rank.  It is
not an LU-orbit boundary and, by
Proposition~\ref{thm:fixed-copy-boundary}, cannot be promoted to a
generic coordinate.
